\subsection{Preesquemas algebraicos geométricamente reducidos.}
\label{subsection:IV.2.4.6-fr}

\begin{proposition}[4.6.1]
\label{IV.2.4.6.1-fr}
Sean $k$ un cuerpo, $X$ un $k$-preesquema y $\Omega$ una extensión perfecta de
$k$. Las condiciones siguientes son equivalentes~:
\begin{enumerate}
\item[\textup{a)}] Para todo $k$-preesquema reducido $S$, $X\times_k S$ es reducido.
\item[\textup{b)}] Para toda extensión $K$ de $k$, $X\otimes_k K$ es reducido.
\item[\textup{c)}] $X\otimes_k\Omega$ es reducido.
\item[\textup{d)}] Para toda extensión radicial finita $k'$ de $k$, $X\otimes_k k'$ es
reducido.
\item[\textup{e)}] $X$ es reducido, y para toda componente irreducible $X_\alpha$ de
$X$, de punto genérico $x_\alpha$, $\boldsymbol{k}(x_\alpha)$ es una extensión separable de
$k$.
\end{enumerate}
\end{proposition}

\begin{proof}
Es trivial que \emph{a)} $\Rightarrow$ \emph{b)} $\Rightarrow$ \emph{c)}~; como $k'$ puede
considerarse una subextensión de $\Omega$, se ha visto (4.4.1) que \emph{c)} implica
\emph{d)}. Demostremos que \emph{d)} implica \emph{e)}. Tomando $k'=k$, se ve primero que $X$ es
reducido~; para demostrar que $\boldsymbol{k}(x_\alpha)$ es separable sobre $k$, basta
demostrar que $\boldsymbol{k}(x_\alpha)\otimes_k k'$ es un anillo reducido para toda extensión
finita radicial $k'$ de $k$ (4.3.5). Es posible limitarse al caso en que
$X=\operatorname{Spec}(A)$ es afín, siendo $A$ una $k$-álgebra reducida~;
entonces, por hipótesis, el anillo $A\otimes_k k'$ es reducido (\textbf I, 5.1.4), y
$\boldsymbol{k}(x_\alpha)\otimes_k k'$ es un anillo de fracciones de este anillo, por lo que
es reducido $(\mathbf{0}_{\mathrm I},1.2.8)$. Finalmente, para demostrar que \emph{e)} implica
\emph{a)}, es posible limitarse al caso en que $X=\operatorname{Spec}(A)$ y
$S=\operatorname{Spec}(B)$ son afines, siendo $A$ y $B$ $k$-álgebras. Los
$\boldsymbol{k}(x_\alpha)$ son los cuerpos de fracciones de los anillos cociente
$A/\mathfrak p_\alpha$, donde $\mathfrak p_\alpha=\mathfrak j_{x_\alpha}$ son los
ideales primos minimales de $A$~; como, por hipótesis,
$\bigcap_\alpha\mathfrak p_\alpha=0$, $A$ está contenido en el producto
$\prod_\alpha A/\mathfrak p_\alpha$ y, por tanto, en
$\prod_\alpha \boldsymbol{k}(x_\alpha)$~; asimismo, puede considerarse $B$ como una
subálgebra de un producto $\prod_\beta L_\beta$ de extensiones de $k$. Se deduce
que $A\otimes_k B$ se identifica con una subálgebra de
$\left(\prod_\alpha \boldsymbol{k}(x_\alpha)\right)\otimes_k
\left(\prod_\beta L_\beta\right)$,
y este producto tensorial se identifica, a su vez, con una subálgebra de
$\prod_{\alpha,\beta}\bigl(\boldsymbol{k}(x_\alpha)\otimes_k L_\beta\bigr)$
(Bourbaki, \emph{Alg.}, chap.~II, 3\textsuperscript{e} éd., \S~7,
n\textsuperscript{o}~7, prop.~15). Ahora bien, por hipótesis, los
$\boldsymbol{k}(x_\alpha)\otimes_k L_\beta$ son reducidos (4.3.5), y lo mismo ocurre, por
tanto, con $A\otimes_k B$.
\end{proof}

\begin{definition}[4.6.2]
\label{IV.2.4.6.2-fr}
Si se satisfacen las condiciones equivalentes \textup{a)} a \textup{e)} de la proposición
\textup{(4.6.1)}, se dice que $X$ es \emph{separable} (o \emph{geométricamente reducido},
o \emph{universalmente reducido}) \textup{sobre} $k$. Se dice que un $k$-preesquema $X$
es \emph{geométricamente íntegro sobre} $k$ si, para toda extensión $K$ de $k$,
$X\otimes_k K$ es íntegro~; equivale (en virtud de \textup{(4.6.1)}) a decir que $X$ es
separable y geométricamente irreducible sobre $k$.
\end{definition}

Se dirá que una $k$-álgebra $A$ (conmutativa) es \emph{separable} si
$\operatorname{Spec}(A)$ es separable sobre $k$~; esto significa, por tanto,
que para toda extensión $K$ de $k$, el anillo $A_{(K)}=A\otimes_k K$ es
reducido. Se observará que esta definición coincide con la de Bourbaki,
\emph{Alg.}, chap.~VIII,

\oldpage[IV-2]{69}

\S~7, n\textsuperscript{o}~5, déf.~1, cuando $A$ es de rango \emph{finito} sobre
$k$ (\emph{loc. cit.}, cor. de la prop.~7), pero no en general, pues una
$k$-álgebra puede tener un radical $\neq0$ incluso si es íntegra.

\begin{corollary}[4.6.3]
\label{IV.2.4.6.3-fr}
Sea $X$ un $k$-preesquema íntegro~; para que $X$ sea geométricamente reducido
\textup{(resp.} geométricamente íntegro\textup{)} sobre $k$, es necesario y suficiente que su
cuerpo de funciones racionales $\mathrm R(X)$ sea una extensión separable
\textup{(resp.} separable y primaria\textup{)} de $k$.
\end{corollary}

\begin{proof}
Esto se deduce inmediatamente de (4.5.9) y (4.6.1).
\end{proof}

\begin{corollary}[4.6.4]
\label{IV.2.4.6.4-fr}
Sea $X$ un $k$-preesquema reducido. Entonces, para toda extensión separable $k'$
de $k$, $X'=X\otimes_k k'$ es reducido.
\end{corollary}

\begin{proof}
Basta aplicar la equivalencia de \emph{a)} y \emph{e)} en (4.6.1), reemplazando $X$ por
$\operatorname{Spec}(k')$ y $S$ por $X$.
\end{proof}

\begin{proposition}[4.6.5]
\label{IV.2.4.6.5-fr}
\begin{enumerate}
\item[\textup{(i)}] Sean $X$ un $k$-preesquema y $K$ una extensión de $k$. Para que $X$
sea geométricamente reducido \textup{(resp.} geométricamente íntegro\textup{)} sobre $k$, es
necesario y suficiente que $X\otimes_k K$ sea geométricamente reducido \textup{(resp.}
geométricamente íntegro\textup{)} sobre $K$.
\item[\textup{(ii)}] Sean $X$, $Y$ dos $k$-preesquemas. Si $X$ e $Y$ son geométricamente
reducidos \textup{(resp.} geométricamente íntegros\textup{)} sobre $k$, lo mismo ocurre con
$X\times_k Y$.
\end{enumerate}
\end{proposition}

\begin{proof}
La afirmación (i) es consecuencia trivial de las definiciones y de (4.4.1).
Para demostrar la parte de la afirmación de (ii) relativa a la separabilidad,
se observa que, si $\Omega$ es una extensión algebraicamente cerrada de $k$, se
tiene
$(X\times_k Y)\otimes_k\Omega=X\times_k(Y\otimes_k\Omega)$
como $Y\otimes_k\Omega$ es reducido, lo mismo ocurre con
$X\times_k(Y\otimes_k\Omega)$ en virtud de (4.6.1, \emph{a)}). El resto de la
afirmación (ii) se deduce de lo anterior y de (4.5.8, (ii)).
\end{proof}

\begin{proposition}[4.6.6]
\label{IV.2.4.6.6-fr}
Sea $X$ un $k$-preesquema de tipo finito. Existe una extensión radicial finita
$k'$ de $k$ tal que $(X_{(k')})_{\mathrm{red}}$ es geométricamente reducido
sobre $k'$.
\end{proposition}

\begin{proof}
Como $X$ puede recubrirse por un número finito de abiertos afines $U_i$, es posible
limitarse al caso en que $X$ es afín~: en efecto, si para cada $i$, $k_i$ es
una extensión radicial finita de $k$ tal que
$(U_i\otimes_k k_i)_{\mathrm{red}}$ es geométricamente reducido sobre $k_i$,
se puede suponer que los $k_i$ están contenidos en una misma extensión radicial
finita $k'$ de $k$~; $(U_i\otimes_k k')_{\mathrm{red}}$ se identifica con
$(U_i\otimes_k k_i)_{\mathrm{red}}\otimes_{k_i}k'$
en virtud de (\textbf I, 5.1.8) y, por tanto, es geométricamente reducido sobre $k'$ por
hipótesis; lo mismo ocurre, pues, con $(X\otimes_k k')_{\mathrm{red}}$.
Supongamos entonces que $X=\operatorname{Spec}(A)$, donde $A$ es una $k$-álgebra
de tipo finito, y sea $\Omega$ una clausura algebraica de $k$. Designemos por
$\mathfrak N'$ el nilradical de $A\otimes_k\Omega$~; como $A\otimes_k\Omega$ es
un anillo noetheriano, $\mathfrak N'$ es un ideal generado por un número finito
de elementos de la forma
$y_i=\sum_j x_{ij}\otimes\xi_{ij}$ donde $x_{ij}\in A$, $\xi_{ij}\in\Omega$.
Sean $K$ una subextensión finita de $\Omega$ que contenga los $\xi_{ij}$, y
$\mathfrak N$ el ideal de $A\otimes_k K$ generado por los $y_i$
($A\otimes_k K$ identificado con un subanillo de $A\otimes_k\Omega$)~; es claro
que los $y_i$ son nilpotentes en $A\otimes_k K$~; por otra parte, como
$\mathfrak N\otimes_K\Omega=\mathfrak N'$ (y, por consiguiente,
$\mathfrak N'\cap(A\otimes_k K)=\mathfrak N$), $\mathfrak N$ contiene el
nilradical de $A\otimes_k K$ y, por tanto, es igual a dicho nilradical. Como
$(A\otimes_k\Omega)/\mathfrak N'
=\bigl((A\otimes_k K)/\mathfrak N\bigr)\otimes_K\Omega$,
se ve que $(X\otimes_k K)_{\mathrm{red}}\otimes_K\Omega$ es reducido y, por
consiguiente, en virtud de (4.6.1), $(X_{(K)})_{\mathrm{red}}$ es
geométricamente reducido sobre $K$. Reemplazando $K$ por la extensión
cuasigaloisiana sobre $k$ que este genera en $\Omega$, se ve por (\textbf I, 5.1.8) que
puede suponerse además que $K$ es
cuasigaloisiana sobre $k$ y, por tanto, extensión separable de una extensión
radicial finita $k'$ de $k$ (Bourbaki, \emph{Alg.}, chap.~V, \S~10,
n\textsuperscript{o}~9, prop.~14). En virtud

\oldpage[IV-2]{70}

de (\textbf I, 5.1.8), $(X_{(k')})_{\mathrm{red}}\otimes_{k'}K$ es isomorfo a
$(X_{(K)})_{\mathrm{red}}$ y, por tanto, es geométricamente reducido~; lo mismo
ocurre entonces con $(X_{(k')})_{\mathrm{red}}$ en virtud de (4.6.5, (i)).
\end{proof}

\begin{corollary}[4.6.7]
\label{IV.2.4.6.7-fr}
Si $K$ es una extensión de tipo finito de $k$, existe una extensión radicial
finita $k'$ de $k$ tal que los cuerpos residuales del anillo semi-local
$K\otimes_k k'$ sean separables sobre $k'$.
\end{corollary}

\begin{proof}
Basta aplicar (4.6.6) a $X=\operatorname{Spec}(A)$, donde $A$ es una
$k$-álgebra de tipo finito cuyo cuerpo de fracciones es $K$.
\end{proof}

\begin{corollary}[4.6.8]
\label{IV.2.4.6.8-fr}
Sea $X$ un $k$-preesquema de tipo finito. Existe una extensión finita $k'$ de
$k$ tal que $(X_{(k')})_{\mathrm{red}}$ es geométricamente reducido sobre $k'$,
que las componentes irreducibles de $X_{(k')}$ son geométricamente irreducibles
y que las componentes conexas de $X_{(k')}$ son geométricamente conexas.
\end{corollary}

\begin{proof}
Esto se deduce inmediatamente de (4.6.6), de (4.5.10) y (4.5.15).
\end{proof}
